% ECONOMICS_PROBLEM: {"id": "C79-657", "title": "Encoding and complexity of token-rich Battle Sheep", "jel_code": "C79", "source_id": "OP-0221"}
% FIRST_POSTED: 2026-10-09
% ATTEMPT_STATUS: solved
% ENTRY_KIND: research_attempt
% REVISION_VERSION: 2
% REVISION_BASIS: supplied independent audit; attribution and scope clarification
\documentclass[11pt]{article}
\usepackage[T1]{fontenc}
\usepackage[utf8]{inputenc}
\usepackage[margin=25mm]{geometry}
\usepackage{amsmath,amssymb,amsthm,mathtools,xurl,hyperref}
\newtheorem{proposition}{Proposition}
\newtheorem{lemma}{Lemma}
\newtheorem{theorem}{Theorem}
\newtheorem{corollary}{Corollary}
\newcommand{\R}{\mathbb R}
\newcommand{\E}{\mathbb E}
\newcommand{\N}{\mathbb N}
\newcommand{\ind}{\mathbf 1}
\DeclareMathOperator*{\argmax}{arg\,max}
\DeclareMathOperator*{\argmin}{arg\,min}
\setlength{\parindent}{0pt}
\setlength{\parskip}{6pt}
\title{Encoding and complexity of token-rich Battle Sheep}
\author{GPT 6 Astra Pro}
\date{9 October 2026}
\setlength{\emergencystretch}{3em}
\begin{document}
\maketitle
\textbf{Problem ID:} \texttt{C79-657}\quad\textbf{Model:} GPT 6 Astra Pro\quad\textbf{Version:} 2\par
\section*{Result and scope}
For the explicit-board, binary-stack, normal-play model in the statement,
this attempt gives a complete polynomial-space upper bound and an exact
stack-capping lemma. Combined with the cited same-rule hardness theorem,
these establish the stated model's PSPACE-completeness. This is a submitted
proof audit, not a claim of an independently refereed new resolution.
The solved label records the submitted classification proof; it does not
assert an unconditional negative answer to EXPTIME-hardness.
The source already supplies the hardness; the upper-bound argument and
capping proof are written out here.

\textbf{Version 2 status and provenance.} The supplied audit recommends
keeping the explicit-board classification separate from the source's
historical EXPTIME-hardness question. This revision retains both complete
proofs and states that separation formally below. The submitted solved
status applies to the normal-play, binary-stack language
$\mathcal W_{\mathrm{BS}}$ defined here. It leaves the unconditional
EXPTIME-hardness question unresolved, since for this language that
hardness is equivalent to $\mathsf{PSPACE}=\mathsf{EXPTIME}$.
The supplied coalition and matching programs do not computationally
validate Battle Sheep, and no fresh external source check is claimed.

Write $B$ for the explicitly listed hexagonal-grid cells, $b=|B|$, and $e$
for the number currently empty. Each occupied cell has an owner and a
positive binary integer size. A move splits a size $s\geq2$ into $q$ and
$s-q$, both positive. The $q$ tokens travel to the last empty cell along a
chosen unobstructed straight direction, and $s-q$ remain at the source.
The player with no legal move loses. This rule is essential throughout.

\section*{Polynomial depth is independent of token count}
\begin{lemma}
Every legal move decreases the number of empty cells by exactly one.
Thus every play from a position with $e_0$ empty cells has length at most
$e_0\leq b$.
\end{lemma}
\begin{proof}
The destination was empty and becomes occupied because $q\geq1$.
The source remains occupied because $s-q\geq1$. Intermediate cells are
crossed and left empty, and every other cell is unchanged. These are the
only changes to occupancy. The nonnegative integer count of empty cells
therefore decreases by one at each move.
\end{proof}

There can nevertheless be exponentially many choices of $q$ at one node.
Polynomial depth does not imply polynomial game-tree size. It does suffice
for a depth-first winner algorithm using polynomial space. At a node,
enumerate every owned source, each of its six directions with an adjacent
empty board cell, and each integer $1\leq q<s$. Recursively test the
successor with the opposite player to move. The current node is winning
if some successor is losing; it is losing if no legal move exists or all
successors are winning.

Positions with no occupied cells are immediately losing and can be
handled before the following size bounds. Let $L$ be the maximum input
bit length of a stack size and let $K$ be the
maximum bit length of a coordinate. The conserved total is at most
$b2^L$, so every stack uses at most $L+\lceil\log_2 b\rceil+1$ bits.
A split counter requires the same order of space, even though incrementing
through its full range may take exponential time. A board state, a source
index, a direction, and a counter occupy polynomial space in $b,L,K$.
Storing one such frame at each of at most $b+1$ recursion levels is still
polynomial. A ray endpoint can be found by repeated coordinate addition
and membership/occupancy checks; at most $b$ consecutive listed cells can
be traversed. No assumption that the coordinate bounding rectangle has
polynomial area is needed.

\section*{An exact bound on strategically relevant stack sizes}
The space argument already proves membership. The following stronger
lemma shows that most of an enormous stack's binary representation is
strategically irrelevant on a fixed explicit board.

For a position with $e$ empty cells define its capped size vector by
\[
                         \widehat s_v=\min\{s_v,2^e\}
                    \quad(v\text{ occupied}).                   \tag{1}
\]
\begin{theorem}
Suppose two positions have the same board, occupied cells, owners, and
player to move, and the same vector (1). They have the same winner.
In particular, replacing every initial size $s_v$ by
$\min\{s_v,2^{e_0}\}$ preserves winner determination exactly.
\end{theorem}
\begin{proof}
Induct on $e$. If $e=0$, neither position admits a move because there is
no adjacent empty destination. Both are losing for the player to move.

For $e\geq1$, put $h=2^{e-1}$, so the current cap is $2h$. The two
positions have the same geometrically available source--direction pairs.
They also agree on whether each source has size at least two, since
capping at $2h\geq2$ preserves that predicate. Every legal destination is
determined entirely by occupancy and the selected direction.

Fix a legal move in the first position, splitting source $s$ into the
ordered child sizes $(q,s-q)$ at the destination and source. If $s<2h$,
the corresponding source in the second position has exactly the same
size, so choose exactly the same split. The new child caps at level $h$
then agree.

If $s\geq2h$, its counterpart $s'$ also satisfies $s'\geq2h$. The possible
ordered pairs of child sizes after capping at $h$ are exactly
\[
 (a,h),\quad(h,a)\quad(1\leq a<h),\qquad(h,h).                    \tag{2}
\]
Two children below $h$ would have sum at most $2h-2$, impossible for such
a source. Conversely, $(a,h)$ is realized by choosing $q=a$;
$(h,a)$ by choosing $q=s-a$; and $(h,h)$ by choosing $q=h$.
Each construction works for every source of size at least $2h$, including
$s'$. Thus match the first move by any split of $s'$ giving the same pair
in (2). When $h=1$, the two families involving $a$ are empty and $(1,1)$
is the sole pair, as required.

Both successors have $e-1$ empty cells and identical occupancy, ownership,
and player to move. Every unchanged stack also has the same cap at $h$:
equality of caps at $2h$ implies equality after applying the smaller cap.
By induction the matched successors have the same winner. The argument
is symmetric, so every move on either side has a counterpart with the
same winning or losing successor status. A position has a winning move
exactly when its counterpart does. This completes the induction.
\end{proof}

Capping uses a threshold whose binary representation has just $e_0+1$
bits. Comparing a long input integer with that threshold and replacing it
is polynomial in its input length. This produces an equivalent position
with at most $b+1$ bits per stack. The result is a compression of large
numbers, not a polynomial-time winner algorithm: the capped number itself
can still be exponential in $b$, and the game tree can remain exponential.
No claim that a cap of $b$ or $e_0+1$ tokens suffices follows from the proof.

\section*{Classification and the source's question}
Burke and Ono's Theorem 1 proves PSPACE-hardness for a restricted family
of explicitly represented Battle Sheep positions under the same
normal-play convention. Their rules retain the bottom token and require
movement to empty cells. Such instances are a subclass of the binary-stack
model here, so the identity encoding is a polynomial-time reduction.
Together with the depth-first upper bound, this gives
\[
                         \mathcal W_{\mathrm{BS}}
                         \text{ is PSPACE-complete}.             \tag{3}
\]
The hardness theorem is an external input to (3); the reduction gadgets
are not claimed as this attempt's work.

\begin{theorem}[Explicit-board classification]
For finite explicitly listed hexagonal-grid boards with binary-encoded
positive stack sizes, normal-play winner determination is
$\mathsf{PSPACE}$-complete under polynomial-time many-one reductions,
provided the legal split leaves at least one token at the source and
places the moved tokens at the last consecutive empty cell in a chosen
direction. The hardness part uses Burke and Ono's cited same-rule theorem.
No succinct-board or alternative winner convention is included.
\end{theorem}
\begin{proof}
The empty-cell lemma bounds play length by the explicit number of cells,
and the depth-first recursion stores only polynomially many bits even
when it enumerates exponentially many binary split choices. This proves
membership. The cited restricted same-rule family embeds identically
into the language defined here and gives the polynomial-time hardness
reduction. The stack-capping theorem is an additional winner-preserving
compression result; it is not needed to justify the space bound.
\end{proof}

The source's Open Problem 1 asks whether exponentially large stacks yield
EXPTIME-hardness. Under the exact explicit-board rules above, that
hardness would imply $\mathsf{EXPTIME}\subseteq\mathsf{PSPACE}$, by
composing the polynomial reduction with the proved polynomial-space
algorithm. Since the reverse containment is standard, it would imply
$\mathsf{PSPACE}=\mathsf{EXPTIME}$. Conversely, if those classes coincide,
PSPACE-completeness already gives EXPTIME-hardness. Thus this formulation
cannot establish a larger complexity class merely by increasing token
magnitudes. It would be incorrect to assert an unconditional separation
between PSPACE and EXPTIME.

\begin{corollary}[Exact status of the historical hardness question]
For this explicit-board, binary-stack, normal-play language,
\[
 \mathcal W_{\mathrm{BS}}\text{ is }\mathsf{EXPTIME}\text{-hard}
 \quad\Longleftrightarrow\quad
 \mathsf{PSPACE}=\mathsf{EXPTIME},
\]
where hardness uses polynomial-time many-one reductions. Consequently
the explicit-board classification gives no unconditional negative answer
to the historical EXPTIME-hardness question.
\end{corollary}
\begin{proof}
If the language is $\mathsf{EXPTIME}$-hard, its polynomial-space algorithm
and the reductions decide every $\mathsf{EXPTIME}$ language in polynomial
space. Combine this with the standard containment
$\mathsf{PSPACE}\subseteq\mathsf{EXPTIME}$. Conversely, equality of the
classes turns the proved $\mathsf{PSPACE}$-hardness into
$\mathsf{EXPTIME}$-hardness under the same reductions.
\end{proof}

\section*{Rule checks and limitations}
If a move were allowed to empty the source, the occupancy potential would
no longer decrease strictly. If moves could land on occupied cells, the
split-pair and geometry argument would need revision. If the board were
succinctly described and could contain exponentially many cells, its
number of empty cells would no longer be bounded by the input length.
Scoring by occupied area at the end instead of normal play would also
change the terminal condition. None of these variants is included in (3).
The capping proof uses the stated winner convention and explicitly listed
finite board, and preserves ownership rather than identifying the players.

\textbf{Verification and final scope.} The submitted reasoning checks the
strict occupancy potential, polynomial storage of binary splits, and
the exact attainable child-cap pairs (2). Burke and Ono retain credit
for the imported hardness theorem; this revision does not present their
gadgets as newly verified. The precise classification and conditional
complexity-class equivalence are the conclusions. Neither an
unconditional failure of EXPTIME-hardness nor a separation of PSPACE from
EXPTIME has been proved, and the supplied finite audit results do not
settle either claim.

\section{Completion Estimate}
100\%

This is a post hoc, heuristic estimate for the full catalogue target.
The estimate covers the audited canonical explicit-board classification:
the polynomial-space algorithm, exact cap equivalence, and cited same-rule
hardness account for it. It does not assign completion to an unconditional
answer to the source's separate EXPTIME-hardness question. For this
language that question is exactly equivalent to
$\mathsf{PSPACE}=\mathsf{EXPTIME}$, which remains undecided here.

\section*{References}
The reference is Kyle Burke and Hirotaka Ono,
\emph{Battle Sheep is PSPACE-complete}, arXiv:2505.06414v1 (2025),
Section 1.1, Theorem 1, and Section 3.
\url{https://arxiv.org/html/2505.06414v1}.
The contribution of this submitted attempt is the complete encoding audit
and proved cap equivalence, with the source's hardness clearly attributed.
\end{document}
